\subsection{\texttt{TLRELU}}
\noindent\textbf{Bundle encoding.} UOP entry 25, opcode \texttt{0x1031}. Bundle: \texttt{B.OP + B.ARG + B.IOR + B.IOTI}. \texttt{B.OP.Opcode}=\texttt{0x1031}. \texttt{B.IOTI}: selector=\texttt{111 last}, \texttt{SrcTile0}=\texttt{src}, \texttt{SrcTile1}=\texttt{absent}, \texttt{DstTile}=\texttt{dst}, \texttt{SizeCode}=profile tile size. \texttt{B.IOR} imports scalar operand(s): \texttt{slope}.\par\smallskip
{\small
\paragraph{PTO ISA description.}
Leaky ReLU with a scalar slope.
\paragraph{Example.}
Tile elements [-2, -1, 0, 1, 2] with alpha=0.01 give [-0.02, -0.01, 0, 1, 2].
\par\smallskip
\paragraph{PTO ISA operation.}
For each element \texttt{(i, j)} in the valid region:

\[
\mathrm{dst}_{i,j} = (\mathrm{src}_{i,j} > 0) ? \mathrm{src}_{i,j} : (\mathrm{src}_{i,j} \cdot \mathrm{slope})
\]
}\par\smallskip
\paragraph{Notes.}
The leaky slope alpha is NOT an explicit operand. It is defined by the element format and the implementation profile. The typical alpha is 0.01 (from the original leaky ReLU paper).
\paragraph{Microcode site table.}
{\scriptsize
\begin{longtable}{@{}R{0.055\linewidth}L{0.23\linewidth}L{0.31\linewidth}L{0.22\linewidth}@{}}
\toprule
Seq & Micro instruction & WO[63:32] fields & WM[31:0] fields \\
\midrule
0 & \texttt{u\_vec.vlds} & \texttt{opc=0x17, x=0, fl=mem, seq=0, kind=b} & \texttt{entry=25, cnt=2, res=vec, var=0} \\
1 & \texttt{u\_vec.vlrelu} & \texttt{opc=0x19, x=0, fl=alu, seq=1, kind=u} & \texttt{entry=25, cnt=1, res=vec, var=0} \\
2 & \texttt{u\_vec.vsts} & \texttt{opc=0x2A, x=0, fl=mem, seq=2, kind=b} & \texttt{entry=25, cnt=2, res=vec, var=0} \\
\bottomrule
\end{longtable}
}
\paragraph{Microcode body DSL.}
\begin{lstlisting}[style=ptocode]
def uop_tlrelu(src, slope, dst):
    dtype = dst.element_type
    valid_rows, valid_cols = dst.valid_shape

    lanes = pto.get_lanes(dtype)
    if pto.constexpr(dtype == pto.f16):
        slope_scalar = pto.f16(slope)
    else:
        slope_scalar = slope

    for row in range(0, valid_rows, 1):
        remained = valid_cols
        for col in range(0, valid_cols, lanes):
            mask, remained = pto.make_mask(dtype, remained)
            src_vec = pto.vlds(src[row, col:])
            result = pto.vlrelu(src_vec, slope_scalar, mask)
            pto.vsts(result, dst[row, col:], mask)
    return
\end{lstlisting}
\Needspace{0.34\textheight}
